%!TEX root = main.tex

Compilers are among the most fundamental and widely-used
software, based on which programs in high-level programming languages
are translated into low-level code. Well known open source compilers include GCC~\cite{GCC} for C and C++ programs and LLVM~\cite{LLVM}.
Compilers for main-stream programming language like C usually include many compilation passes and optimizations. As a result of their inherent complexity,
there remain massive bugs in compilers, despite intensive efforts of testing~\cite{LeAS14,ZhangSS17}.
Even worse, compiler bugs could silently entail incorrect executions in low-level programs and lead to catastrophic consequences
for safety-critical applications, and it is extremely difficult to determine whether
an execution failure is caused by defects in programs or
bugs in compilers. This motivates researchers to apply formal verification techniques to guarantee the correctness of compilers. 

Formally verified compilers have been built for several realistic programming languages, such as CompCert for C~\cite{Leroy09,kastner2018compcert,Leroy2006,Leroy-BKSPF-2016,BlazyDL06,Kastner-LBSSF-2017},
V\'{e}lus for Lustre~\cite{BourkeBDLPR17}, Vellvm~\cite{ZhaoNMZ12,ZhaoNMZ13} and Crellvm~\cite{KangKSLPSKCCHY18} for LLVM Intermediate Representation,
and CakeML for ML~\cite{KumarMNO14,LoowKTMNAF19,TanMKFON19,OwensNKMT17}.
Such compilers come with a mathematical, machine-checked proof
that compilation passes preserve program semantics,
hence strengthening the correctness guarantees that can be obtained by formal verification of high-level programs.

The success of formally verified compilers on software programming languages has been ported to hardware description languages (HDLs)
very recently, e.g., the extension of CakeML (called Lutsig) for compiling %register-transfer language Verilog to gate-level netlist~\cite{Loow21,loowphd21}
a subset of register-transfer level (RTL) Verilog programs to netlist level programs for field-programmable gate array (FPGA),
and the extension of CompCert (called Vericert) for compiling %compiling C programs to Verilog programs~\cite{HerklotzPRW21}.
a subset of C programs to RTL Verilog programs~\cite{HerklotzPRW21}.
Following this promising research direction, in this work, we establish for the first time a framework that formally reasons about programs in FIRRTL~\cite{firrtl-spec} (Flexible Intermediate Representation for RTL) and its lowering transformations (aka compilation passes).  

FIRRTL originated from the high-level hardware construction language Chisel~\cite{BachrachVRLWAWA12}.
Chisel is a domain specific language embedded in Scala that comprises a library of special class definitions, predefined objects, and usage conventions for writing highly-parameterized circuit design generators.
High-level hardware generators in Chisel are developed by manipulating circuit components using Scala functions,
and their interfaces are encoded in Scala types.
Circuits can also be constructed by leveraging Scala's functional and object-oriented programming features, enabling parameterized, modular, and reusable hardware generators that improve hardware design productivity and robustness.
Chisel has been successfully used to build RISC-V processes such as Rocket Chip~\cite{RocketChip}, RISC-V BOOM \cite{Boom}, NutShell~\cite{NutShell}, 
XiangShan~\cite{xiangshan,Xiangshan-github}, and so on. 
%SiFive Core IP portfolio~\cite{SiFive},
%and Google Edge TPU~\cite{EdgeTPU}.
The FIRRTL intermediate representation was proposed for the compilation of Chisel programs~\cite{IzraelevitzKLLW17}, analogous to the role of LLVM IR~\cite{LLVM} for compiling C programs.
FIRRTL features first-class support for high-level constructs
such as vector types, bundle types, conditional statements, connects, partial connects,
and modules so that Chisel programs can easily be translated into FIRRTL.
%
High-level synthesis is performed on FIRRTL by a sequence of compilation passes and optimizations, resulting into a most restricted form that is close to RTL Verilog, called low FIRRTL (abbreviated as LoFIRRTL).
While FIRRTL was designed to compile Chisel programs, it is not fundamentally tied to Chisel. Other HDLs can also be compiled into FIRRTL and reuse the majority of
the FIRRTL compilation passes. For instance, FIRRTL has been integrated into the CIRCT project~\cite{CIRCT} as one of its core IRs.

Partially due to the fact that Chisel and FIRRTL are relatively new and the ecosystem around them is still evolving, formalization and verification of FIRRTL IR and its lowering transformations are largely missing, with an exception \cite{HarrisonHA21}. In \cite{HarrisonHA21}, a coinductive denotational semantics for LoFIRRTL based on a formalism called device calculus was proposed. \zhilin{a comparison with this work should be carefully done. }

 operational semantics of FIRRTL is still missing and
the FIRRTL compiler suffers from lots of bugs due to complexity of FIRRTL transformations.
In fact,  there are 649 issues in its github repository, among which
99 have been labeled as bugs~\cite{FIRRTLissures}.
For instance, the width of \lstinline!dshl(e1,e2)! inferred by the ``Infer Width'' compilation pass
is less than the width of \lstinline!e1!~\footnote{\url{https://github.com/chipsalliance/firrtl/issues/2296}}, where \lstinline!dshl! is a dynamic shift left operation
which shifts the bits in \lstinline!e1! towards the
most significant bit \lstinline!e2! places and \lstinline!e2! zeroes are shifted in to the least significant bits.
The ``ExpandWhens'' compilation pass may introduce unexpected \lstinline!VOID!, leading to a match error in
the ``InferTypes'' compilation pass~\footnote{\url{https://github.com/chipsalliance/firrtl/issues/852}}.

We define for the first time the executable operational semantics of low FIRRTL.
present the first formal operational semantics of the hardware description language FIRRTL (Flexible Intermediate Representation for RTL)~\cite{firrtl-spec} and verify two key compilation passes of FIRRTL programs~\cite{firrtlcompiler,IzraelevitzKLLW17}, namely ``Infer Width'' and ``Expand Whens'',
which infers bit widths of signals
and replaces conditional statements with multiplexers, respectively.





\fu{In this work, we make the first attempt to fill the above gap. We tbd. Describe what we have done, e.g., challenges of formalizing
operational semantics of FIRRTL and transformations, how we address these challenges,
results of our work.}
